\documentclass[10pt,a4paper]{amsart}
\usepackage[margin=19mm]{geometry}
\usepackage{amsmath,amssymb,mathtools}
\usepackage[scr=rsfs,cal=euler]{mathalfa}
\usepackage{xfrac}
\usepackage[unicode,hidelinks,bookmarks=false]{hyperref}
\newcommand{\ie}{i.e.\ }
\newcommand{\cf}{cf.\ }
\newcommand{\resp}{respectively\ }
\newcommand{\Subsection}{\subsection}
\providecommand{\nobreakdash}{\nobreak\hbox{-}}
\setlength{\emergencystretch}{2em}
\multlinegap0pt
\usepackage{bm}
\usepackage{bbm}
\usepackage{dsfont}
\usepackage{xspace}
\usepackage{cancel}
\usepackage{mathtools}
\usepackage{amsthm}
\makeatletter
\@ifundefined{theorem}{\newtheorem{theorem}{Theorem}}{}
\@ifundefined{lemma}{\newtheorem{lemma}[theorem]{Lemma}}{}
\makeatother
\newcommand{\Z}{\mathbb{Z}}
\title[Cyclotomic norm congruences for zeta values and modular $L$-]{Cyclotomic norm congruences for zeta values and modular $L$-values}
\author{Shubhrajit Bhattacharya, Anwesh Ray, R. Sujatha}
\date{Source: arXiv:2608.22738v1 (CC BY 4.0)}
\begin{document}
\maketitle

\noindent\emph{Source and licence.} Cyclotomic norm congruences for zeta values and modular $L$-values. Version arXiv:2608.22738v1 (CC BY 4.0), distributed under the Creative Commons Attribution 4.0 International licence, as recorded in the arXiv licence metadata for this exact version and shown on the abstract page https://arxiv.org/abs/2608.22738v1. Full rights notice: licenses/arxiv-2608.22738-CC-BY-4.0.txt.\par\smallskip

\noindent\textit{Editorial scope.} Counted unit: the Lemma whose source label is \texttt{lem:A1-A0} in the article (source0.tex lines 3571--3585, source label \texttt{lem:A1-A0}), delivered together with the article's own complete original proof of it (source0.tex lines 3587--3977, one \texttt{proof} environment of 391 lines). No mathematics is written, repaired or re-derived: statement and proof are the article's own text line for line. Required context retained uncounted: lem:Ani-integral-formula (lines 2769--2799); eq:A0i-interpolation (lines 3176--3184); eq:trivial-algebraic-integrality (lines 3197--3201); eq:A00-bottom (lines 3405--3411); eq:p2-inflated-gauss-sum (lines 3440--3448); ramsum (lines 3544--3548); eq:ramanujan-p2 (lines 3554--3563). Every definition, display and label of those lines is retained unchanged. Omitted from this package and not redistributed: 1--2768, 2800--3175, 3185--3196, 3202--3404, 3412--3439, 3449--3543, 3549--3553, 3564--3570, 3586--3586, 3978--4220. Nothing inside a retained excerpt is omitted or altered: every retained body is delivered exactly as the article's lines are written, so no omission receipt of any kind accompanies this package. Citations: the retained excerpts carry no citation command at all, so no bibliography entry of the article is transcribed and no reference is redistributed. Reference adaptation: none was needed and none was made; every reference inside the delivered unit resolves inside it, so no unresolved reference or citation is produced. Printed slips kept literal: no printed word, symbol, hypothesis or number of the counted statement or of its proof was altered. Layout and class: the article's own class is \texttt{amsart} with packages including geometry, amscd, amsmath, amsxtra, amsthm, amssymb, stmaryrd, xr, mathrsfs, mathtools, enumerate, commath, comment, orcidlink; the wrapper is the lane's pinned \texttt{amsart} editorial wrapper at 10pt on A4, which restates the theorem-like environments the retained text needs and adds the article's own macro definitions that the retained text needs (\texttt{\textbackslash Z}), with extra wrapper packages (bm, bbm, dsfont, xspace, cancel, mathtools). The delivered numbering is the unit's own. Assets: no retained span contains a figure environment, a graphics-inclusion command, a PDF-page inclusion, a table or a TikZ picture; no image file is copied into this package, the package declares no asset dependency and no image file is redistributed with it. Licence: version arXiv:2608.22738v1 of this article is distributed under the Creative Commons Attribution 4.0 International licence, as recorded in the arXiv licence metadata for this exact version and shown on the abstract page https://arxiv.org/abs/2608.22738v1. A full rights notice is delivered at licenses/arxiv-2608.22738-CC-BY-4.0.txt. Characters: every retained excerpt is pure ASCII, so the delivered engine, XeLaTeX with its default Latin Modern text font, reports no missing character.
\begin{lemma}
\label{lem:Ani-integral-formula}
Let \(n\geq0\) and \(0\leq i\leq p-2\), and put
\[
        \varepsilon_i:=\omega^i(-1)\in\{\pm1\}.
\]
Then
\begin{equation}
\label{eq:Ani-modular-symbol-formula}
        A_{n,i}
        =
        \sum_{a\in(\Z/p^{n+1}\Z)^\times}
        \omega^i(a\bmod p)\,
        \varphi_f^{\varepsilon_i}
        \left(
        \{\infty\}
        -
        \left\{\frac{a}{p^{n+1}}\right\}
        \right)\Big|_{(X,Y)=(0,1)}.
\end{equation}
Equivalently,
\begin{equation}
\label{eq:Ani-integral-formula}
        A_{n,i}
        =
        \frac{2\pi i}{\Omega_f^{\varepsilon_i}}
        \sum_{a\in(\Z/p^{n+1}\Z)^\times}
        \omega^i(a\bmod p)
        \int_{a/p^{n+1}}^{i\infty}f(z)\,dz.
\end{equation}
\end{lemma}

\begin{equation}
\label{eq:A0i-interpolation}
        A_{0,i}
        =
        \tau(\omega^i)
        \frac{\mathrm{L}(f,\overline{\omega^i},1)}
             {\Omega_f^{\omega^i(-1)}}
        \in\mathcal O.
\end{equation}

\begin{equation}
\label{eq:trivial-algebraic-integrality}
        \frac{\mathrm{L}(f,1)}{\Omega_f^+}
        \in\mathcal O.
\end{equation}

\begin{equation}
\label{eq:A00-bottom}
        A_{0,0}
        =
        \bigl(1-a_p+p^{k-2}\bigr)
        \frac{\mathrm{L}(f,1)}{\Omega_f^+}.
\end{equation}

\begin{equation}
\label{eq:p2-inflated-gauss-sum}
        S_i(r)
        =
        \begin{cases}
        0,&p\nmid r,\\[3pt]
        p\,\tau(\omega^i)\overline{\omega^i}(t),&r=pt.
        \end{cases}
\end{equation}

\begin{equation}\label{ramsum}c_q(r)
        :=
        \sum_{a\in(\Z/q\Z)^\times}
        e^{2\pi i ar/q}.
\end{equation}

\begin{equation}
\label{eq:ramanujan-p2}
        c_{p^2}(r)
        =
        \begin{cases}
        0,&p\nmid r,\\
        -p,&p\mid r,\ p^2\nmid r,\\
        p(p-1),&p^2\mid r.
        \end{cases}
\end{equation}

\begin{lemma}
\label{lem:A1-A0}
For every \(0\leq i\leq p-2\),
\begin{equation}
\label{eq:A1-A0}
        A_{1,i}
        \equiv
        a_pA_{0,i}
        \pmod{\mathfrak m}.
\end{equation}
For \(i\neq0\), one in fact has the exact identity
\[
        A_{1,i}=a_pA_{0,i}.
\]
\end{lemma}

\begin{proof}
We first treat the case \(1\leq i\leq p-2\). Set $\chi:=\omega^i$ and $\varepsilon:=\chi(-1)$. The character \(\chi\) is primitive of conductor \(p\). By Lemma \ref{lem:Ani-integral-formula},
\begin{equation}
\label{eq:A1i-integral}
        A_{1,i}
        =
        \frac{2\pi i}{\Omega_f^\varepsilon}
        \sum_{a\in(\Z/p^2\Z)^\times}
        \chi(a)
        \int_{a/p^2}^{i\infty}f(z)\,dz.
\end{equation}

We now evaluate the integrals in \eqref{eq:A1i-integral} using the
Fourier expansion of \(f\).  Since the points \(a/p^2\) lie on the
boundary of the upper half-plane, we first move them a distance
\(y>0\) into the upper half-plane.  For \(y>0\),
\[
        2\pi i
        \int_{a/p^2+iy}^{i\infty}f(z)\,dz
        =
        -\sum_{r\geq1}
        \frac{a_r(f)}{r}
        e^{2\pi i ar/p^2}e^{-2\pi ry}.
\]
The series on the right is absolutely convergent for every \(y>0\).
Since the sum over \(a\in(\Z/p^2\Z)^\times\) is finite, we may
therefore interchange the sum over \(a\) with the Fourier series.
Taking the limit \(y\to0^+\) only after this rearrangement gives
\[
\begin{aligned}
        A_{1,i}
        &=
        -\frac1{\Omega_f^\varepsilon}
        \lim_{y\to0^+}
        \sum_{a\in(\Z/p^2\Z)^\times}
        \chi(a)
        \sum_{r\geq1}
        \frac{a_r(f)}{r}
        e^{2\pi i ar/p^2}e^{-2\pi ry}             \\
        &=
        -\frac1{\Omega_f^\varepsilon}
        \lim_{y\to0^+}
        \sum_{r\geq1}
        \frac{a_r(f)}{r}e^{-2\pi ry}
        \sum_{a\in(\Z/p^2\Z)^\times}
        \chi(a)e^{2\pi i ar/p^2}.
\end{aligned}
\]

The factor \(p\) in \eqref{eq:p2-inflated-gauss-sum} cancels the factor
\(p\) in the denominator \(r=pt\) of the Fourier integral. Thus, with
the same modular-symbol and Gauss-sum conventions as in
\eqref{eq:A0i-interpolation}, the Abel-regularized Fourier expression
for \(A_{1,i}\) is
\[
        A_{1,i}
        =
        \frac{\tau(\chi)}{\Omega_f^\varepsilon}
        \lim_{y\to0^+}
        \sum_{t\geq1}
        \frac{a_{pt}(f)}{t}
        \overline{\chi}(t)e^{-2\pi pty}.
\]
On the other hand, \eqref{eq:A0i-interpolation} gives
\[
        A_{0,i}
        =
        \frac{\tau(\chi)}{\Omega_f^\varepsilon}
        \lim_{y\to0^+}
        \sum_{t\geq1}
        \frac{a_t(f)}{t}
        \overline{\chi}(t)e^{-2\pi ty}.
\]
We now use the Hecke relation at \(p\nmid N\):
\[
        a_{pt}(f)
        =
        a_pa_t(f)-p^{k-1}a_{t/p}(f),
\]
where \(a_{t/p}(f)=0\) if \(p\nmid t\). Substituting this into the
expression for \(A_{1,i}\), we obtain
\[
\begin{aligned}
        A_{1,i}
        &=
        \frac{\tau(\chi)}{\Omega_f^\varepsilon}
        \lim_{y\to0^+}
        \left(
        a_p
        \sum_{t\geq1}
        \frac{a_t(f)}{t}
        \overline{\chi}(t)e^{-2\pi pty}
        \right.\\
        &\hspace{45mm}\left.
        -
        p^{k-1}
        \sum_{t\geq1}
        \frac{a_{t/p}(f)}{t}
        \overline{\chi}(t)e^{-2\pi pty}
        \right).
\end{aligned}
\]
The second sum vanishes term by term. Indeed, if \(p\nmid t\), then
\(a_{t/p}(f)=0\), whereas if \(p\mid t\), then
\(\overline{\chi}(t)=0\). Hence
\[
        A_{1,i}
        =
        a_p\frac{\tau(\chi)}{\Omega_f^\varepsilon}
        \lim_{y\to0^+}
        \sum_{t\geq1}
        \frac{a_t(f)}{t}
        \overline{\chi}(t)e^{-2\pi pty}.
\]
Since \(py\to0\) as \(y\to0^+\), the last Abel limit is
\(\mathrm{L}(f,\overline{\chi},1)\). Consequently,
\[
        A_{1,i}
        =
        a_p\tau(\chi)
        \frac{\mathrm{L}(f,\overline{\chi},1)}
             {\Omega_f^\varepsilon}
        =
        a_pA_{0,i},
\]
where the last equality is \eqref{eq:A0i-interpolation}. This proves
the exact identity
\[
        A_{1,i}=a_pA_{0,i}
        \qquad
        (1\leq i\leq p-2).
\]
It remains to treat the case when \(i=0\). By Lemma \ref{lem:Ani-integral-formula}, we find that
\[
        A_{1,0}
        =
        \frac{2\pi i}{\Omega_f^+}
        \sum_{a\in(\Z/p^2\Z)^\times}
        \int_{a/p^2}^{i\infty}f(z)\,dz.
\]
To justify the Fourier expansion at the lower endpoints, we again use
Abel regularization.  For \(y>0\),
\[
        2\pi i
        \int_{a/p^2+iy}^{i\infty}f(z)\,dz
        =
        -\sum_{r\geq1}
        \frac{a_r(f)}{r}
        e^{2\pi i ar/p^2}e^{-2\pi ry},
\]
and the series on the right is absolutely convergent.  Hence, for
fixed \(y>0\), we may sum over \(a\) and interchange the finite sum
with the Fourier series.  It follows that
\[
\begin{aligned}
        A_{1,0}
        &=
        -\frac1{\Omega_f^+}
        \lim_{y\to0^+}
        \sum_{a\in(\Z/p^2\Z)^\times}
        \sum_{r\geq1}
        \frac{a_r(f)}{r}
        e^{2\pi i ar/p^2}e^{-2\pi ry} \\
        &=
        -\frac1{\Omega_f^+}
        \lim_{y\to0^+}
        \sum_{r\geq1}
        \frac{a_r(f)}{r}e^{-2\pi ry}
        \sum_{a\in(\Z/p^2\Z)^\times}
        e^{2\pi i ar/p^2}.
\end{aligned}
\]
By the definition of the Ramanujan sum \eqref{ramsum}, the inner sum is
\(c_{p^2}(r)\). Therefore
\[
        A_{1,0}
        =
        -\frac1{\Omega_f^+}
        \lim_{y\to0^+}
        \sum_{r\geq1}
        \frac{a_r(f)}r e^{-2\pi ry}c_{p^2}(r).
\]
Using \eqref{eq:ramanujan-p2}, the coefficient of a term with
\(p\mid r\), \(p^2\nmid r\), is \(p\), while the coefficient of a term
with \(p^2\mid r\) is \(-p(p-1)\). These two cases can be combined as
\[
        A_{1,0}
        =
        \frac1{\Omega_f^+}
        \lim_{y\to0^+}
        \left(
        p\sum_{p\mid r}
        \frac{a_r(f)}r e^{-2\pi ry}
        -
        p^2\sum_{p^2\mid r}
        \frac{a_r(f)}r e^{-2\pi ry}
        \right).
\]
Indeed, for \(p^2\mid r\) the coefficient on the right is
\(p-p^2=-p(p-1)\), as required. We evaluate the two terms separately. In the first sum, write \(r=pm\).
Then
\[
\begin{aligned}
        p\sum_{p\mid r}
        \frac{a_r(f)}r e^{-2\pi ry}
        &=
        \sum_{m\geq1}
        \frac{a_{pm}(f)}m e^{-2\pi pmy}.
\end{aligned}
\]
Using
\[
        a_{pm}(f)
        =
        a_pa_m(f)-p^{k-1}a_{m/p}(f),
\]
we obtain
\[
\begin{aligned}
        \sum_{m\geq1}
        \frac{a_{pm}(f)}m e^{-2\pi pmy}
        &=
        a_p
        \sum_{m\geq1}
        \frac{a_m(f)}m e^{-2\pi pmy}        \\
        &\quad
        -
        p^{k-1}
        \sum_{m\geq1}
        \frac{a_{m/p}(f)}m e^{-2\pi pmy}.
\end{aligned}
\]
The first sum is \(F(py)\). In the second sum only multiples \(m=pn\)
contribute, so
\[
\begin{aligned}
        p^{k-1}
        \sum_{m\geq1}
        \frac{a_{m/p}(f)}m e^{-2\pi pmy}
        &=
        p^{k-1}
        \sum_{n\geq1}
        \frac{a_n(f)}{pn}e^{-2\pi p^2ny} \\
        &=
        p^{k-2}F(p^2y).
\end{aligned}
\]
Therefore
\begin{equation}
\label{eq:first-p2-Abel-sum}
        \sum_{m\geq1}
        \frac{a_{pm}(f)}m e^{-2\pi pmy}
        =
        a_pF(py)-p^{k-2}F(p^2y).
\end{equation}
Since \(F(u)\to\mathrm{L}(f,1)\) as \(u\to0^+\), it follows that
\begin{equation}
\label{eq:first-p2-Abel-limit}
        \lim_{y\to0^+}
        \sum_{m\geq1}
        \frac{a_{pm}(f)}m e^{-2\pi pmy}
        =
        (a_p-p^{k-2})\mathrm{L}(f,1).
\end{equation}
For the second term, writing \(r=p^2m\) gives
\[
        p^2\sum_{p^2\mid r}
        \frac{a_r(f)}r e^{-2\pi ry}
        =
        \sum_{m\geq1}
        \frac{a_{p^2m}(f)}m e^{-2\pi p^2my}.
\]
Applying the Hecke relation to the index \(pm\), we have
\[
        a_{p^2m}(f)
        =
        a_pa_{pm}(f)-p^{k-1}a_m(f).
\]
Consequently,
\[
\begin{aligned}
        \sum_{m\geq1}
        \frac{a_{p^2m}(f)}m e^{-2\pi p^2my}
        &=
        a_p
        \sum_{m\geq1}
        \frac{a_{pm}(f)}m e^{-2\pi p^2my}       \\
        &\quad
        -
        p^{k-1}
        \sum_{m\geq1}
        \frac{a_m(f)}m e^{-2\pi p^2my}.
\end{aligned}
\]
The second sum is \(F(p^2y)\). Applying
\eqref{eq:first-p2-Abel-sum} with \(y\) replaced by \(py\) to the
first sum gives
\[
        \sum_{m\geq1}
        \frac{a_{pm}(f)}m e^{-2\pi p^2my}
        =
        a_pF(p^2y)-p^{k-2}F(p^3y).
\]
Hence
\[
\begin{aligned}
        \sum_{m\geq1}
        \frac{a_{p^2m}(f)}m e^{-2\pi p^2my}
        &=
        a_p
        \bigl(
        a_pF(p^2y)-p^{k-2}F(p^3y)
        \bigr)
        -
        p^{k-1}F(p^2y).
\end{aligned}
\]
Letting \(y\to0^+\), we find
\begin{equation}
\label{eq:second-p2-Abel-limit}
\begin{aligned}
        \lim_{y\to0^+}
        \sum_{m\geq1}
        \frac{a_{p^2m}(f)}m e^{-2\pi p^2my}
        &=
        \left(
        a_p(a_p-p^{k-2})-p^{k-1}
        \right)\mathrm{L}(f,1).
\end{aligned}
\end{equation}
Combining \eqref{eq:first-p2-Abel-limit} and
\eqref{eq:second-p2-Abel-limit}, we obtain
\[
\begin{aligned}
        A_{1,0}
        &=
        \frac1{\Omega_f^+}
        \left[
        (a_p-p^{k-2})
        -
        \bigl(
        a_p(a_p-p^{k-2})-p^{k-1}
        \bigr)
        \right]
        \mathrm{L}(f,1)                                      \\
        &=
        \frac1{\Omega_f^+}
        \left(
        a_p-a_p^2+a_pp^{k-2}-p^{k-2}+p^{k-1}
        \right)
        \mathrm{L}(f,1)                                      \\
        &=
        \left[
        a_p(1-a_p+p^{k-2})
        +
        p^{k-2}(p-1)
        \right]
        \frac{\mathrm{L}(f,1)}{\Omega_f^+}.
\end{aligned}
\]
By \eqref{eq:A00-bottom},
\[
        a_pA_{0,0}
        =
        a_p(1-a_p+p^{k-2})
        \frac{\mathrm{L}(f,1)}{\Omega_f^+}.
\]
Subtracting, we therefore obtain the exact identity
\[
        A_{1,0}-a_pA_{0,0}
        =
        p^{k-2}(p-1)
        \frac{\mathrm{L}(f,1)}{\Omega_f^+}.
\]
Since \(k>2\), we have \(p^{k-2}\in\mathfrak m\), while
\(\mathrm{L}(f,1)/\Omega_f^+\in\mathcal O\) by
\eqref{eq:trivial-algebraic-integrality}. Hence
\[
        A_{1,0}-a_pA_{0,0}\in\mathfrak m,
\]
or equivalently
\[
        A_{1,0}
        \equiv
        a_pA_{0,0}
        \pmod{\mathfrak m}.
\]
Together with the exact identity
\(A_{1,i}=a_pA_{0,i}\) for \(1\leq i\leq p-2\), this proves
\eqref{eq:A1-A0}.
\end{proof}

\end{document}
